\subsection{MULTIPLIABLE SEQUENCES IN A NORMED ALGEBRA}

Let A be a normed algebra over a non-discrete valued field K (Chapter IX, § 3, no. 7, Definition 9); we shall denote by \(||x||\) the norm of an element \(x \in A\) , and we shall assume that this norm satisfies the inequality \(||xy|| \leqslant ||x|| \cdot ||y||\) ; also we shall assume that A has an identity element \(e\) . Let \((x_n)_{n \in \mathbb{N}}\) be an infinite sequence of points of A. Every finite subset J of N, linearly ordered by the ordering of N, defines a sequence \((x_n)_{n \in J}\) of points of A, and we define the product

\[
p _ {J} = \prod_ {n \in J} x _ {n}
\]

of this sequence; this product is called the finite partial product of the sequence \((\pmb{x}_n)_{n\in \mathbb{N}}\) , corresponding to the finite subset \(J\) of \(\mathbb{N}\) (recall that if \(J = \emptyset\) we put \(\prod_{n\in \emptyset}x_n = e\) ).

DEFINITION 1. The sequence \((\pmb{x}_n)_{n\in \mathbb{N}}\) is said to be multipliable in the normed algebra \(\mathbf{A}\) if the mapping \(\mathbf{J}\to \pmb{p}_{\mathbf{J}}\) has a limit with respect to the filter of sections of the set \(\mathfrak{F}(\mathbf{N})\) of finite subsets of \(\mathbf{N}\) , ordered by the relation \(c\) ; this limit is called the product of the sequence \((\pmb{x}_n)_{n\in \mathbb{N}}\) , and is denoted by \(\prod_{n\in \mathbb{N}}\pmb{x}_n\) (or simply \(\prod_{n}\pmb{x}_n\) ); the \(\pmb{x}_n\) are called the factors of this product.

Definition 1 is equivalent to the following: the sequence \((\pmb{x}_n)\) is multipliable and its product is \(\pmb{p}\) if for each \(\varepsilon > 0\) there exists a finite subset \(J_0\) of \(\mathbf{N}\) such that, for every finite subset \(J \supset J_0\) of \(\mathbf{N}\) , we have \(||\pmb{p}_{\mathrm{J}} - \pmb{p}|| \leqslant \varepsilon\) .

Remarks. 1) When A is a commutative algebra, Definition 1 is identical with that given in Chapter III, § 5, no. 1, Remark 3; but when A is not commutative, the order structure of the index set N is essentially involved in Definition 1. If \(\sigma\) is an arbitrary permutation of N, we cannot in general assert that the sequence \((\pmb{x}_{\sigma(n)})\) is multipliable if the sequence \((x_{n})\) is multipliable; and if both sequences are multipliable, their products will in general be different.

\begin{enumerate}
\def\labelenumi{\arabic{enumi})}
\setcounter{enumi}{1}
\tightlist
\item
  Definition 1 can be immediately generalized to the case of a family \((\pmb{x}_n)_{n\in I}\) whose index set I is a subset of \(\mathbf{Z}\) (linearly ordered by the order induced by that of \(\mathbf{Z}\) ). We leave it to the reader to extend to this case the results below (cf.~Exercises 1 and 2).
\end{enumerate}

